\chapter{基础算法及技巧}

\section{基础算法}

\subsection{排列前缀的 MEX}

对 $0,1,\ldots,n-1$ 的排列建立后缀最小值数组。查询前 $k$ 项的 MEX 时取 \texttt{suffixMin[k]}，整个数组的 MEX 为 $n$。

\begin{pseudocode}
suffixMin[n] = n
for i = n - 1 .. 0:
    suffixMin[i] = min(a[i], suffixMin[i+1])
前 k 项的 MEX = suffixMin[k]
\end{pseudocode}

\section{搜索算法}

\subsection{折半搜索与二进制枚举}

适用于整集枚举超时、两半状态可以按条件配对的题目。$n=40$ 时，每半枚举 $2^{20}$ 个掩码。

\begin{enumerate}
    \item 按下标将元素拆成大小相差至多一的两组。对每组独立枚举二进制掩码，记录每个子集参与合并的状态，例如元素和或选取个数。
    \item 根据题目要求的关系组织右侧状态。若要求左右状态和为 $T$，左侧状态 $x$ 只需查询右侧的 $T-x$。等值查找可以用哈希或排序后二分，区间查找可分别求上下界。
    \item 统计方案数时保留重复状态的次数。若同时限制总共选 $k$ 项，左侧选了 $i$ 项时，只需查询右侧选 $k-i$ 项的状态。
\end{enumerate}

填写空子集的 \texttt{initial}，并让 \texttt{take(old,i)} 返回追加第 \texttt{i} 项后的新状态。每半的元素数须小于掩码位宽。

\cppfile{code/basic/meet_in_the_middle.cpp}

右侧排序后，对每个左侧状态二分查询。枚举存储 $2^{\lfloor n/2\rfloor}+2^{\lceil n/2\rceil}$ 个状态。

\section{位运算技巧}

\subsection{二进制位的读写与集合运算}

从最低位起按 $0$ 编号。令 $m=2^k$，下列操作只涉及第 $k$ 位。移位量须小于无符号类型的位宽。

\begin{cppcode}
// x 为 ull，0 <= k < 64。
ull mask = 1ULL << k;
bool bit = (x & mask) != 0;
ull setBit = x | mask;
ull clearBit = x & ~mask;
ull flipBit = x ^ mask;
\end{cppcode}

把置位的位置看作集合，按位与、按位或、按位异或分别对应交集、并集、对称差。
若 $x$ 是全集掩码 $U$ 的子掩码，其补集为 $U\oplus x$。
掩码 $y$ 是 $x$ 的子集当且仅当 $(x\mathbin{\mathtt{\&}}y)=y$。

\subsection{异或恒等式与进位}

求只出现一次的值且其余值均出现偶数次时，将所有元素异或。需要拆分加法或按位或时可直接用下式。
\[
a+b=(a\oplus b)+2(a\mathbin{\mathtt{\&}}b),\qquad
a\mathbin{\mathtt{|}}b=(a\oplus b)+(a\mathbin{\mathtt{\&}}b).
\]
按普通整数计算时须选用足够宽的类型。

\subsection{最低置位}

对于非零无符号数 $x$，$x\mathbin{\mathtt{\&}}(-x)$ 保留最低置位，$x\mathbin{\mathtt{\&}}(x-1)$ 清除最低置位。前一式在 $x=0$ 时返回零。正整数是二的幂当且仅当只有一个置位。

\begin{cppcode}
// x 为 ull。无符号取负按 2^64 取模。
ull low = x & -x;
bool isPowerOfTwo = x != 0 && (x & (x - 1)) == 0;
int count = 0;
for (ull s = x; s; s &= s - 1) ++count;
\end{cppcode}

循环次数等于置位数。直接求置位数、尾零与前导零的内建函数见第八章。

\subsection{二进制枚举}

令 \texttt{n} 为元素个数。用 \texttt{mask} 处理每个子集，内层循环给出该子集的元素下标。

\cppfile{code/basic/binary_enumeration.cpp}

只枚举掩码耗时 $O(2^n)$。对每个掩码再遍历其置位时，$n\ge 1$ 下共访问 $n2^{n-1}$ 个元素。

\subsection{子掩码与超掩码}

传入全集掩码 \texttt{mask}，在循环内处理每个子掩码 \texttt{s}，包含空集。

\cppfile{code/basic/submask.cpp}

原掩码有 $k$ 个置位时，共枚举 $2^k$ 个子掩码。
遍历所有 $n$ 位掩码及各自子掩码的总次数为 $3^n$。

若要求枚举 $x$ 在全集 $U$ 内的超掩码，先枚举 $U\oplus x$ 的子掩码 $s$，再取 $x\mathbin{\mathtt{|}}s$。

\subsection{枚举恰有 $k$ 个置位的掩码}

指定总位数 \texttt{n} 和选取数量 \texttt{k}，在循环内处理当前子集。掩码按数值递增，$k=0$ 时只处理空集。

\cppfile{code/basic/fixed_popcount.cpp}

设有效位数为 $n$。在 $0\le k\le n<64$ 下，生成所有掩码的时间为 $O\bigl(\binom nk\bigr)$，额外空间为 $O(1)$。


\subsection{连续整数的按位异或}

按下面的方式计算前缀异或 $S(n)$。查询非负闭区间 $[l,r]$ 时取 $S(r)\oplus S(l-1)$。

\begin{pseudocode}
if n < 0: return 0
按 n % 4 为 0、1、2、3，分别返回 n、1、n+1、0
\end{pseudocode}
